\section{FT的实际实现}

section 2描述了FT的基本设计和协议。然而，要创建一个可用、健壮且自动化的系统，%
还必须设计并实现许多其他组件。

\subsection{启动和重启FT VM}

必须设计的最大额外组件之一，是一种能够以与primary VM%
相同状态启动backup VM的机制。这个机制也会在发生failure%
之后重新启动backup VM时使用。因此，该机制必须能够用于一个%
处于任意状态的正在运行的primary VM(也就是说，并不只是用于%
启动阶段)。此外，我们希望该机制不会显著干扰primary VM的执%
行，因为这会影响该VM当前的任何客户端。%

对于VMware FT，我们改造了VMware vSphere现有的VMotion功%
能。VMware VMotion[10]允许把一个正在运行的VM从一台服务器迁移到%
另一台服务器，并且造成的干扰很小，VM暂停时间通常小于一秒。%
我们创建了一种修改后的VMotion，它会在远程服务器上创建一个VM的精确运行副%
本，但不会销毁本地服务器上的VM。也就是说，我们修改后的FT VMotion会把一%
个VM克隆到远程host，而不是迁移它。FT VMotion还会建立logging %
channel，并使源VM以primary VM身份进入recording mode%
，使目标VM以新backup VM身份进入replay mode。%
与普通VMotion一样，FT VMotion通常使primary VM的执行暂停%
不到一秒。因此，在一个正在运行的VM上启用FT是一个简单且非破%
坏性的操作。%
\\\emph{译者注：
老式操作系统可能主要依赖timer interrupt来推进系统时间。例如，%
硬件每10ms产生一次定时器中断，那么内核收到10次中断后，就认为大约%
过去了100ms。现代操作系统通常不会仅靠这种周期性tick来计算当前时%
间，而是会选择某个clocksource，例如x86机器上的TSC%
(Time Stamp Counter，CPU时间戳计数器)，或者虚拟化环境中的%
paravirtual clock。调用gettimeofday时，系统会读取当前%
clocksource的计数值，与内核timekeeper中保存的基准计数值%
(如cycle\_last)相减，将cycle差值换算成纳秒，再加到内核维%
的wallclock基准时间上，最后返回给用户态。很多系统还会通过%
vDSO在用户态完成这一步，避免每次都陷入内核。
\\针对VM暂停和迁移，读者可能会问：在复制期间，如果VM大约暂停了1s，%
是否会导致VM内部读到的时间落后现实世界时间(wallclock)1s？可%
以这样理解：如果系统真的完全依赖timer interrupt计数来推进时间，%
并且暂停期间没有收到这些中断，那么恢复后确实可能出现时间落后；但现%
代虚拟化环境通常不会简单依赖补发timer interrupt来追赶时间。%
hypervisor可以通过虚拟化TSC、虚拟时钟偏移、paravirtual clock%
或时间同步机制，让guest OS在下一次读取clocksource或调用%
gettimeofday时看到已经校正后的wallclock(real world clock)。%
因此，VM暂停1s后，不需要向guest OS连续注入大量timer interrupt来弥补这1s。%
wallclock的推进和timer interrupt的注入是两个不同层面的问题；VM恢复时，%
hypervisor会根据虚拟定时器的类型和语义处理已经过期的timer，可能合并或投递%
timer interrupt，并不保证固定补发一个中断。
\\\\读取wallclock也是一种non-deterministic operation。原因是wallclock
来自VM外部，如果primary VM和backup VM分别读取各自host的真实时间，
它们可能得到不同结果，从而导致执行路径分叉。VM FT的做法是：primary VM
执行这类读取操作时，由hypervisor返回一个虚拟时间值，并把该返回值及其发生位置
写入replay log；backup VM replay到同一位置时，不读取自己的真实时间，而是直接
使用log中记录的值。这样，primary和backup看到的时间完全一致。
需要注意的是，用户态的gettimeofday不一定每次都会直接产生log entry。
如果它只是读取guest memory中由guest OS维护的时间变量，那么这个普通内存读取
本身是确定性的。真正需要记录的是把外部时间引入VM的边界操作，例如读取虚拟时钟
设备、执行RDTSC，或者timer interrupt推进guest OS时间状态。}%

启动backup VM的另一个方面，是选择在哪台服务器上运行它。%
fault-tolerant VM运行在能够访问shared storage的服务器cluster中，%
因此所有VM通常都可以在cluster中的任意服务器上运行。%
这种灵活性使VMware vSphere即使在一台或多台服务器已经failure时%
，也能恢复FT冗余。VMware vSphere实现了一个维护管理和资源信息的%
cluster服务。当发生failure并且某个primary VM现在需要一个新%
的backup VM来重新建立冗余时，primary VM会通知clustering service%
它需要一个新的backup。clustering service根据资源使用%
情况和其他约束确定运行backup VM的最佳服务器，并调用FT VMotion来创%
建新的backup VM。结果是，VMware FT通常能够在服务器failure后%
的几分钟内重新建立VM冗余，并且整个过程不会对fault-tolerant VM的%
执行造成任何可察觉的interrupt。

\subsection{管理logging channel}

\insertfigurethree

管理logging channel上的流量涉及一些有趣的实现细节。在我们的实现中，%
hypervisor为primary VM和backup VM维护一个大的log %
entry缓冲区。当primary VM执行时，它会把log entry生成到log %
buffer中；类似地，backup VM会从自己的log buffer消费log %
entry。primary VM log buffer的内容会尽快刷新到logging %
channel，log entry也会在到达后尽快从logging channel读%
入backup VM的log buffer。backup VM每次从网络读入一些log %
entry到自己的log buffer时，都会向primary VM发送%
acknowledgment。这些acknowledgment使VMware FT%
能够确定那些被\textbf{Output Rule} delay的输出何时可以发送。%
Figure 3展示了这一过程。\\\emph{译者注：log entry包含，%
timer/disk-write-completion interrupt、wallclock read、%
鼠标/键盘/network/disk-read-completion input，以及为了满足%
\textbf{Output Rule}的output log entry，这些output包括%
写共享磁盘、写外部设备，还有对外写网络报文。}%

如果backup VM在需要读取下一条log entry时遇到空的log buffer%
，它会停止执行，直到有新的log entry可用\emph{(译者注：%
backup已经没有下一条%
log entry可读，因此它不知道primary接下来是否会在某个精确指%
令位置收到interrupt、disk completion、timer interrupt%
等非确定性事件。即使backup接下来执行的普通指令本身是确定性的，%
它也不能擅自继续执行，因为下一条非确定性事件可能本应发生在下一条%
或很近的某条指令之后。backup只有在读到下一条log entry后，才知%
道可以安全执行到哪个指令位置，并在该位置注入相同的非确定性事件。%
因此，backup不是“看到确定性代码就能一直跑”，而是“在已知下一条%
log entry边界的前提下，可以连续执行边界之前的确定性指令”。%
当log buffer为空时，这个边界未知，所以必须暂停)}。%
由于backup VM不与外部通信%
，这种暂停不会影响VM的任何客户端。类似地，如果primary VM在需要写入%
log entry时遇到满的log buffer，它必须停止执行，%
直到log entry能够被刷新出去。这种执行停止是一种自然的流控机制，%
会在primary VM过快产生log entry时降低其速度。然而，%
这种暂停可能影响VM的客户端，因为primary VM会完全停止并且无响应，%
直到它能够record log entry并继续执行。因此，%
我们的实现必须被设计成尽量降低primary log buffer被填满的可能性。%
\\\emph{译者注：简化的primary执行与日志生成伪代码；实际实现还包含%
并发的日志发送、ACK处理和buffer回收逻辑：}
\begin{Verbatim}[
  fontsize=\scriptsize,
  breaklines=true,
  breakanywhere=true
]
// VM执行循环
func PrimaryRunLoop(vm *VM) {
    for vm.Running {
        result := RunGuestUntilExitOrEvent(vm)

        switch result.Type {

        case "DETERMINISTIC_EXECUTION":
            // 普通CPU指令、内存读写、guest OS内部逻辑
            // 不写log，backup自己也能执行出来
            continue

        case "NONDETERMINISTIC_EVENT":
            // 例如timer interrupt、network input、disk I/O completion
            entry := LogEntry{
                Seq:        NextSeq(),
                EventType:  result.EventType,
                InstrCount: vm.InstrCount(),
                Payload:    result.Payload,
            }

            // 关键：primary把非确定性事件写入自己的logbuffer
            primaryLogBuffer.Append(entry)

            // 然后才把这个事件注入/交给guest
            InjectEventToGuest(vm, entry)

        case "OUTPUT_OPERATION":
            // guest想对外输出，比如发网络包
            // 注意：output本身可能是deterministic产生的，
            // 但它有外部副作用，所以要遵守OutputRule

            entry := LogEntry{
                Seq:        NextSeq(),
                EventType:  "OUTPUT_MARKER",
                InstrCount: vm.InstrCount(),
                Payload:    nil,
            }

            primaryLogBuffer.Append(entry)

            output := PendingOutput{
                Packet:      result.Packet,
                RequiredSeq: entry.Seq,
            }

            // 不能马上发出去，先放到pending队列
            pendingOutputs.Append(output)

        case "VM_FAIL_OR_STOP":
            return
        }
    }
}
// 一个单独的线程，作为背景线程运行
func HandleBackupAck(ackedSeq uint64) {
    primaryLogBuffer.ReleaseThrough(ackedSeq)

    for pendingOutputs.HasReady(ackedSeq) {
        output := pendingOutputs.PopReady(ackedSeq)
        SendToExternalWorld(output.Packet)
    }
}
\end{Verbatim}
\emph{如果Append因buffer已满而无法继续，primary执行循环就必须等待。%
日志发送循环把entry发送给backup；收到ACK后，ACK处理逻辑回收已确认的%
log entry，并释放满足Output Rule的pending output。}

primary log buffer可能被填满的一个原因，是backup VM执行得太慢，%
因而消费log entry也太慢。一般来说，backup VM必须能够以大致等于%
primary VMrecord执行的速度来replay执行。幸运的是，%
在VMware deterministic replay中，%
record和replay的开销大致相同。然而，如果承载backup VM的服务器%
上运行了许多其他VM(因而资源过度提交)，那么即使backup hypervisor的VM%
调度器尽力调度，backup VM也可能无法获得足够的CPU和内存资源来以%
primary VM的速度执行。
\\\emph{译者注：这里的内存资源不足，并不是说guest OS看到%
的内存容量变小了。创建VM时确实会指定一个内存大小，例如4GB%
或8GB，guest OS也会认为自己拥有这么多“物理内存”。但是这%
只是guest看到的guest physical memory大小，这些guest物理页背%
后仍然需要宿主机(host)真实物理内存承载。如果承载backup VM的机器上%
运行了很多VM，所有VM配置的内存总量可能超过物理机的真实物理内存，%
这就是memory overcommit，也就是内存超卖。此时hypervisor可能%
通过ballooning、memory compression、host swapping、page sharing%
等机制回收或腾挪内存页。如果host物理内存实在不够，hypervisor%
也可能从其他VM回收分配给该VM的物理页；直观地说，这有点像向其他VM%
“借”内存，但不是把另一台VM整体永久暂停后直接搬走它的内存。%
更准确地说，hypervisor可以把其他VM中暂时不用、可回收或者%
优先级较低的页面压缩后放入host内存中的压缩区，或者写入host swap所在的%
磁盘空间，也可以通过ballooning%
让该guest OS主动释放一些页面，再把释放出来的host物理页分配%
给当前更需要资源的VM。被回收页面的VM通常仍然可以继续运行，%
但如果它之后再次访问这些被压缩或换出的页面，就需要等待hypervisor%
重新解压或从swap读回，因此会发生短暂停顿。更底层地说，guest OS%
维护自己的页表，把guest virtual address映射到guest physical address；hypervisor%
还会维护第二级地址映射，例如Intel EPT或AMD NPT，把guest physical%
address映射到host physical address。当某个guest物理页暂时没有映射到host%
物理页时，guest再次访问该页不会由guest OS自己处理，而是触发%
EPT violation或nested page fault，使CPU退出到hypervisor。随后%
hypervisor把该页从swap中读回，或者解压、重新分配host物理页，%
并恢复对应的二级页表映射。这个过程更准确地说是VM-exit，不是%
普通外部中断。它不会改变guest OS看到的内存大小，但会增加内存%
访问延迟。因此backup VM虽然名义上仍拥有创建时配置的内存容量，%
却可能因为host层面的物理内存竞争、换页、向其他VM回收页面，%
或者NUMA局部性变差而replay变慢。}%

除了避免log buffer被填满导致意外暂停之外，我们不希望执行滞后变得太大还有%
另一个原因。如果primary VM failure，backup VM必须通过%
replay它已经acknowledged的所有log entry来“追上来”%
，然后才能go live并开始与外部世界通信。完成replay所需时间基本上就是%
failure发生点的执行滞后时间，因此backup VM go live所需时间大致%
等于failure detection时间加上当前执行滞后时间。因此，%
我们不希望执行滞后时间很大(超过一秒)，因为这会显著增加failover%
(故障接管)的时间。\emph{(译者注：因为backup必须重放完，直到处理完%
最后一条log entry)}%

因此，我们有一个额外机制来降低primary VM的速度，%
防止backup VM落后太多。在发送和acknowledgment log entry%
的协议中，我们会发送额外信息，以确定primary VM和%
backup VM之间的实时执行滞后。通常执行滞后小于100毫秒。%
如果backup VM开始出现显著执行滞后(例如超过1秒)，%
VMware FT会通过通知调度器给primary VM略少一些CPU(开始时只减少%
几个百分点)来减慢primary VM。我们使用一个缓慢的反馈循环，%
逐步找出适合primary VM的CPU限制，使backup VM能够匹配其执行。%
如果backup VM继续落后，我们会继续逐步降低primary VM的CPU限制。%
相反，如果backup VM追上来了，我们会逐步提高primary VM的CPU限制%
，直到backup VM恢复到只有轻微滞后。%
\\\emph{译者注：\textbf{Output Rule}会把对外输出延迟到backup VM ACK对应%
log entry之后。例如logging channel的一次确认往返若耗时约40ms，额外出包%
延迟可能接近这一量级；实际延迟还取决于log flush、批处理和调度。}%

注意，primary VM的这种减速非常少见，通常只会在系统处于极端压力下发生。%
section 5中的所有性能数字都包含了任何这类减速的成本。

\subsection{FT VM上的操作}

另一个实际问题是处理可能应用于primary VM的各种控制操作。例如，%
如果primary VM被显式关机，backup VM也应停止，%
而不应尝试go live。再如，primary VM上的任何resource management%
变化(例如增加CPU份额)也应应用到backup VM。%
对于这类操作，会通过logging channel从primary VM向%
backup VM发送特殊的control entry，以便在backup VM上执行%
相应操作。%

一般来说，对VM的大多数操作都只应在primary VM上发起。%
VMware FT随后会发送必要的control entry，%
使backup VM发生相应变化。唯一可以在primary VM和backup VM上%
独立执行的操作是VMotion。也就是说，primary VM和backup VM可%
以独立地VMotion到其他host。注意，VMware FT会确保两个VM都不会%
被移动到另一个VM所在的服务器上，因为那种情况将不再提供%
fault tolerance能力。%

primary VM的VMotion比普通VMotion增加了一些复杂性，%
因为backup VM必须在适当时刻断开与源primary VM的连接，%
并重新连接到目标primary VM。backup VM的VMotion也有类似问题%
，但还增加了额外复杂性。对于普通VMotion，我们要求在VMotion最终切换%
发生时，所有未完成的磁盘IO都被静止(也就是完成)。对于primary VM，%
通过等待物理IO完成并把这些完成投递给VM，就很容易处理这种静止。然而，%
对于backup VM，没有简单方法能在任何要求的位置使所有IO都完成，%
因为backup VM必须replay primary VM的执行，%
并在相同的执行点完成IO。primary VM可能正在运行一种workload，%
在正常执行期间始终有磁盘IO正在进行。VMware FT有一种独特方法来解决这个问%
题。当backup VM到达VMotion的最终切换点时，%
它会通过logging channel请求primary VM临时静止其所有IO。%
随后，backup VM在replay primary VM的静止操作执行时，%
它自己的IO也会自然地在同一个执行点被静止。
\\\emph{译者注：这是说了两个问题：一是primary从一台host迁移到另外一台；%
二是backup从一台host迁移到另外一台。其中后者更难一点。
\\\\primary迁移可以理解为：在最终切换点暂停primary VM，等待所有in-flight%
磁盘I/O完成，并把这些I/O completion交付进VM的虚拟设备状态/待处理中断状态中，%
从而形成一个没有inflight disk I/O的干净迁移点。随后旧primary停止运行，%
把旧VM状态切换到新host，并且配置网络，新primary从该状态继续执行。%
backup则在正确的log边界断开旧primary，并连接新primary。%
此后新primary继续作为primary记录non-determinism事件，生成log entry，%
并通过新的logging channel发送给backup。%
\\\\backup迁移比primary迁移更麻烦：原因是backup VM并不是自由执行的VM，%
而是在replay primary VM的执行流；因此backup不能为了迁移而任意选择一个时刻%
交付磁盘I/O completion，否则就可能破坏和primary相同的执行轨迹。%
当backup VMotion进入最终切换阶段时，backup会通过logging channel请求primary%
暂时quiesce磁盘I/O。primary执行这个quiescing operation，等待已有in-flight%
磁盘I/O完成，并把completion交付给VM；该操作随后被记录进FT log。backup在%
replay primary执行流时，会在同一个执行点replay到这个quiescing operation，%
于是backup的I/O也自然进入quiesced状态。此时backup才可以安全地完成迁移，%
旧backup停止运行，目标主机上的backup从该状态继续replay primary后续产生的log。}%

\subsection{磁盘IO的实现问题(重要，难点)}

\emph{译者注：此处我解释磁盘IO工作流程。guest OS启动并初始化磁盘驱动时，会在%
guest物理内存中分配Descriptor Table、Available Ring和Used Ring。%
Descriptor Table用于描述请求涉及的内存buffer；Available Ring由guest驱动%
写入、虚拟设备读取，用于发布新请求；Used Ring由虚拟设备写入、guest驱动%
读取，用于返回已经完成的请求。%
\\Descriptor Table中的一个descriptor包含buffer的guest物理地址、长度、%
访问方向以及下一个descriptor的编号。一个磁盘请求不一定只使用一个%
descriptor，而可能由一条descriptor chain组成。例如，一个write请求通常%
包含请求头、待写数据和完成状态三个部分：请求头和待写数据由guest写入、%
设备读取，完成状态则由设备写入、guest读取。guest驱动会记录descriptor %
chain的head ID，以便之后根据该ID找到原始请求。%
\\Available Ring中保存的不是完整IO请求，而是descriptor chain的head ID。%
avail.idx是位于guest物理内存中的16位计数器，由guest驱动维护，%
并不是一个MMIO寄存器。假设avail.idx初始值为0，guest准备提交%
三个分别由descriptor 7、12和4描述的请求，那么guest会依次执行%
avail.ring[0]=7、avail.ring[1]=12和%
avail.ring[2]=4。填写完descriptor和Available Ring后，guest必须%
执行适当的memory barrier，保证设备在观察到新的avail.idx时，%
前面的descriptor内容已经可见；随后guest把avail.idx更新为3，%
最后写doorbell MMIO寄存器，通知设备有新的请求可供处理。%
\\设备内部会维护一个类似lastAvailIdx的消费位置。收到doorbell后，%
设备比较自己的lastAvailIdx和guest写入的avail.idx，%
然后依次读取尚未消费的Available Ring元素，并根据其中的head ID遍历%
对应的descriptor chain。Available Ring是循环数组，因此实际数组位置为%
逻辑index对queue size取模；但是在设备把请求写入Used Ring以前，guest不能%
修改或复用该请求占用的descriptor及相关buffer。如果队列已经没有空闲%
descriptor，guest驱动必须暂时把新请求保留在自己的软件队列中，等待旧请求%
完成，而不能直接覆盖仍由设备使用的descriptor。%
\\在协议没有规定额外顺序约束的情况下，设备可以并发处理多个请求，%
其物理执行顺序和completion顺序不一定与Available Ring中的提交顺序相同。%
例如，三个请求的提交顺序可能是7、12、4，但实际完成顺序可能是12、4、7。%
如果两个write访问同一个sector，而guest又没有通过等待前一个请求完成、%
FLUSH、FUA、barrier或其他协议机制建立顺序关系，那么设备不保证这两个%
write按照提交顺序生效，最终sector中保存的可能是任意一个write的数据。%
这里更准确的说法是最终结果不受虚拟磁盘协议保证，而不是发生了类似%
C/C++ undefined behavior的任意行为。%
\\某个请求完成时，设备首先完成对descriptor所描述buffer的访问。对于read，%
设备会把读取到的数据写入device-writable data buffer；对于read和write，%
设备还会把成功或失败状态写入对应的status buffer。完成这些写入并执行%
必要的memory barrier后，设备把该请求的descriptor chain head ID写入%
used.ring[used.idx \% queueSize]，随后增加used.idx，%
并根据当前的通知策略向guest OS发送interrupt。%
\\guest磁盘驱动维护一个自己的lastUsed。收到interrupt后，驱动会%
处理从lastUsed到最新used.idx之间的所有Used Ring元素。%
每个Used Ring元素都包含原请求的descriptor chain head ID，因此驱动可以%
通过该ID找到原始block IO请求，检查status，使用read返回的数据，完成对应%
的IO，唤醒等待线程，并回收descriptor和buffer。Used Ring中的数组位置%
不需要与Available Ring中的提交位置一一对应；真正用于匹配请求和completion%
的是descriptor ID。只有在Used Ring表明请求已经完成后，guest才能安全地%
复用该请求占用的descriptor。%
\\一个completion也不一定对应一次interrupt。为了减少interrupt和VM-exit%
开销，设备可以连续完成多个请求、填写多个Used Ring元素，然后只发送一次%
interrupt。guest收到这次interrupt后，会通过used.idx发现并处理%
全部新增的completion。%
\\在虚拟机中，上述设备行为通常由VMM中的虚拟磁盘设备实现。真实host IO%
完成后，VMM把read数据和status写入guest memory，更新Used Ring和%
used.idx，然后向guest注入虚拟磁盘interrupt。对于VMware FT，%
Primary VMM还需要记录IO completion的结果及其在guest指令流中的投递位置；%
Backup VMM正常replay时不会执行真实磁盘IO，而是在收到并replay相应日志后，%
写入相同的数据和status、更新相同的Used Ring，并在相同的guest执行位置%
注入interrupt。}%

与磁盘IO相关的实现问题有许多微妙之处。首先，由于disk operation是非%
阻塞的，因此可以并行执行，同时访问同一disk location的%
并发disk operation可能导致non-determinism。此外，%
我们的磁盘IO实现使用DMA直接读写VM的内存，因此同时访问同一内存页的并发磁盘%
操作也可能导致non-determinism。我们的解决方案通常是检测任何这样的%
IO竞争(这种情况很少见)，并强制这些竞争的disk operation在%
primary VM和backup VM上以相同方式顺序执行。%
\\\emph{译者注：这里讨论的是同一个VM提交的并发磁盘操作。如果多个操作%
访问相同的disk location或guest memory page，不同的执行或完成顺序可能%
产生不同结果；VMware FT需要检测这类I/O竞争，并保证它们在primary和%
backup上按照相同顺序串行执行。这不是primary和backup之间争用shared disk。%
\\典型disk read流程是：驱动在内存中填写I/O descriptor，其中包含磁盘位置、%
长度和目标内存位置，再写doorbell寄存器通知设备。控制器读取磁盘后通过DMA%
写入目标内存，并报告completion。disk write则由控制器通过DMA读取待写数据。%
在默认设计中，disk read返回的数据及其completion是input，会通过log提供给%
backup；disk write是output。非共享磁盘设计会让backup执行disk write，而让%
backup执行disk read是section 4.2讨论的另一项独立设计。}%

其次，disk operation也可能与VM中应用(或OS)的内存访问竞争，%
因为disk operation通过DMA直接访问VM的内存。例如，%
如果VM中的应用/OS正在读取一个内存块，同时又有一个disk read正在读入该%
内存块，就可能产生non-determinism结果。这种情况也不太可能发生，%
但一旦发生，我们必须检测并处理它。一种解决方案是临时对作为%
disk operation目标的页面设置page protection。%
如果VM恰好访问了一个也是未完成disk operation目标的页面，%
page protection会导致一次trap，VM可以被暂停直到%
disk operation完成。由于更改页面上的MMU保护是昂贵操作，%
我们选择改用bounce buffers。bounce buffers是一个临时缓冲%
区，其大小与disk operation访问的内存相同。disk read操作被修改%
为把指定数据读到bounce buffers中，并且只有在IO完成被投递时，%
才把数据复制到guest内存。类似地，对于disk write操作，%
要发送的数据会先复制到bounce buffers，disk write也被修改为从%
bounce buffers写出数据。使用bounce buffers可能减慢%
disk operation，但我们没有看到它造成任何可察觉的性能损失。%
\\\emph{译者注：这里是一个具体的edge case，我们仔细讲。%
\\\\现在有一个primary和backup，primary中的guest OS发起了一次disk read，%
请求把某个磁盘block读入guest内存页A；后续guest CPU又会执行读取%
内存页A的指令。由于磁盘I/O是异步的，如果不加控制，磁盘设备可能会通过%
DMA在某个不可控时刻把数据写入A。在primary上，guest提交读操作后不久直接读内存A，%
由于primary没有等到中断，可能读到的是旧数据，也可能是磁盘读入的数据，%
甚至还可能是混杂数据。而backup在默认设计中不会执行真实disk read；在对应%
completion log entry到达前，它的内存A仍保持旧值。%
这样primary和backup在同一条读取指令上就可能读到不同的值，从而导致VM状态分叉。%
\\\\解决办法是使用bounce buffer。primary发生disk read时，hypervisor不让磁盘设备%
直接DMA写入guest内存页A，而是先把数据DMA到hypervisor的临时缓冲区%
bounce buffer中；当磁盘IO完成时暂停VM，将bounce buffer%
内存数据拷贝进内存A，并向guest注入磁盘完成的interrupt，%
记录completion内容及其精确执行位置并生成log entry，随即恢复VM执行。%
backup在replay时也会遇到guest提交disk read的动作；在默认设计中，backup%
不会把该请求真正发往磁盘，而是等待primary记录的read结果和completion。%
后续读取内存A时，它会读到旧值，%
当停到上面提到的来自primary的disk interrupt log entry后，backup的%
hypervisor才会暂停VM，紧接着在拿着log entry里的信息在内存A里面写数据%
，随即注入interrupt。}%

第三，当failure发生且backup VM执行take over时，%
primary VM上in-flight的磁盘IO会带来一些问题。%
新提升的primary VM无法确定这些磁盘IO是否已经发往磁盘，%
或是否已经成功完成。此外，由于这些磁盘IO没有在backup VM上对外发出，%
随着新提升的primary VM继续运行，也不会有明确的IO完成返回给它们，%
这最终会导致VM中的guest操作系统启动abort或reset过程。%
我们可以发送一个错误完成，表示每个IO都失败了，因为即使IO已经成功完成，%
返回错误也是可接受的。然而，guest OS可能无法很好地处理来自本地磁盘的错误。%
相反，我们在backup VM的go live过程中重新发出这些待处理IO。%
由于我们已经消除了所有竞争，并且所有IO都直接指定了访问哪些内存和%
disk block，因此即使这些disk operation已经成功完成，%
也可以重新发出它们(也就是说，它们是幂等的)。
\\\emph{译者注：hypervisor向外发网络包和hypervisor向共享磁盘发请求，%
表面上都属于VM对外部世界的output，因此都受到\textbf{Output Rule}约束：%
一个output至少要等到backup确认收到足够的log之后，primary才可以真正%
把它释放到外部世界；在此之前，该output只能暂存在hypervisor中。%
\\\\对于网络报文，\textbf{Output Rule}只保证：如果某个packet已经对外可见，那么%
backup一定已经拥有足够的log，可以在failover后进入与该输出一致的状态。%
但是，\textbf{Output Rule}并不保证primary释放packet之后，该packet一定已经完整%
通过物理网卡发出。primary可能在packet尚未真正发出、只发出一部分、或者%
已经完整发出之后发生failure。backup无法区分这些情况，也不会在go live%
之后补发这个packet。原因是网络本身允许丢包、重复和乱序：如果packet没有%
真正发出去，可以视为普通网络丢包；如果packet已经发出去，backup再补发%
反而可能制造重复输出。对于TCP，这类丢包或重复会由序号、ACK和重传机制%
处理；对于UDP，则本来就允许丢包。因此，网络输出的不确定性可以被网络%
协议语义吸收，FT层不需要也不应该猜测并补发旧packet。%
\\\\共享磁盘的磁盘I/O则不同。guest OS把本地磁盘看作块设备，请求发出后通常期待一个%
明确的I/O completion。如果old primary上的pending disk I/O已经发往共享%
存储，但backup原先并没有真正发出这些I/O，那么backup接管后也不会自然%
收到这些completion。若什么都不做，guest OS中的块设备请求会悬挂；若直接%
返回本地磁盘错误，又可能触发abort或reset流程。因此，VMware FT选择在%
backup go live过程中重新发出这些pending disk I/O。这里能够重发，并不是%
因为磁盘I/O天然像网络协议一样支持任意重传，而是因为系统已经消除了I/O%
竞争，并且每个I/O都明确指定了访问的guest memory范围和disk block范围。%
在这些条件下，普通block read/write可以被安全地重新执行：同一份数据重复%
写入同一批block不会改变最终磁盘状态；没有竞争写入时，重新读取同一批%
block也会得到一致结果。具体实现可能是：
\\hypervisor在虚拟磁盘设备层维护一张%
pending disk I/O表。每当guest OS向虚拟磁盘控制器提交一个I/O请求时，%
hypervisor就创建一条PendingDiskIO记录，其中包含请求类型、访问的disk block范围、%
目标guest memory范围、对应的bounce buffer、请求状态以及completion是否%
已经交付给guest等信息。对于write，hypervisor会先把guest memory中要写出的数据%
复制到bounce buffer，再从bounce buffer向共享存储发出真实write；这样即使%
guest随后修改原内存页，这个write请求的数据内容也保持稳定。对于read，%
hypervisor不会让真实磁盘DMA直接写入guest memory，而是先读入bounce buffer，%
等到一个确定的I/O completion投递点时，才把bounce buffer中的数据复制到%
guest memory，并更新虚拟磁盘设备状态、向guest注入完成中断。%
在primary上，hypervisor首先根据guest提交的请求建立PendingDiskIO。%
对于write，hypervisor会立即把guest memory中的待写数据复制到bounce %
buffer，从而固定本次write的数据内容；对于read，则为返回数据准备bounce %
buffer。需要注意，read请求本身不会修改外部状态，因此primary可以把真实read%
提交给共享存储，但是read返回的数据、完成状态和completion投递时机都属于%
non-deterministic input，必须通过FT log复制给backup。%
\\磁盘write则不同。对shared storage的write会改变外部世界，因此属于output，%
必须遵循Output Rule。primary不能在建立PendingDiskIO之后立即提交真实write，%
而是要先在FT log中加入一条与该output关联的特殊log entry，并把该entry以及%
此前所有相关日志发送给backup。只有收到backup的ACK，确认backup已经收到足够%
的日志后，primary才可以把bounce buffer中的数据真正写入shared storage。%
这个ACK只表示backup已经收到日志，不要求backup已经replay到该write的位置；%
即使backup仍存在执行滞后，它也可以在primary故障后继续replay，并最终恢复到%
能够重建该pending write的状态。%
\\物理I/O完成后，primary VMM会先捕获completion状态；对于read，还会保存%
返回到bounce buffer中的数据。随后primary记录completion结果及其在guest%
指令流中的确定投递位置，并生成对应的FT log entry。完成日志记录后，%
primary才在该执行位置把read数据复制到guest memory，更新虚拟磁盘设备状态，%
并向guest注入completion interrupt。只有completion已经交付给guest以后，%
该请求才从primary的pending disk I/O状态中结束。%
\\在backup上，guest提交I/O的动作会通过deterministic replay再次执行，因此%
Backup VMM能够恢复出对应的虚拟磁盘请求和pending状态。正常replay期间，%
Backup不会把这些read或write真正提交给shared storage。对于write，Backup%
收到Output Rule所需的日志时只保存日志并返回ACK，不会因此执行真实write；%
当它replay到primary发送的completion log entry时，才在相同的guest执行位置%
写入相同的read数据和completion状态，更新虚拟设备，并向guest注入相同的%
interrupt。完成交付后，对应请求才从Backup的pending状态中移除。%
\\如果primary发生failure，Backup首先replay完已经收到的FT log，然后经过%
防止split-brain的仲裁过程进入go-live状态。此时它扫描仍然pending的磁盘I/O：%
已经通过completion log交付给guest的请求不再执行；尚未交付completion的请求%
则根据保存的请求状态重新提交给shared storage。对于write，重新提交的是相同%
bounce buffer中的数据和相同disk block；对于read，则重新读取相同的disk %
block。这样，无论旧primary尚未提交I/O、已经提交但尚未完成，还是已经完成但%
completion log尚未到达Backup，新primary都能通过重新执行pending I/O获得一个%
明确的completion。}%

\subsection{网络IO的实现问题}

VMware vSphere为VM网络提供了许多性能优化。%
其中一些优化基于hypervisor异步更新VM网络设备的状态。例如，%
hypervisor可以在VM执行时直接更新接收缓冲区。不幸的是，%
这些对VM状态的异步更新引入了non-determinism。%
除非我们能保证所有更新都发生在primary VM和backup VM指令流中的同一%
位置，否则backup VM的执行就可能偏离primary VM。%

为了支持FT，对网络仿真代码做出的最大改动，是禁用异步网络优化。%
原本会用传入包异步更新VM ring buffer的代码，已经被修改为强制guest %
trap到hypervisor；在hypervisor中，%
系统可以记录这些更新，然后再把它们应用到VM。类似地，%
通常会从transmit queue异步取出包的代码在FT中也被禁用，%
发送改为通过一次到hypervisor的trap完成(下面注明的例外除外)。%
\\\emph{译者注：
\\
这里的问题与disk read completion类似，核心都在于虚拟设备可能异步修改
guest-visible state。普通vSphere为了提高网络性能，可能会在guest VM
继续执行时，由hypervisor异步更新receive buffer、descriptor ring或
transmit queue等网络设备状态。但这些更新如果没有被记录到replay log中，
就可能在primary VM和backup VM的指令流中发生于不同位置，从而破坏
deterministic replay。
\\因此，在FT模式下，VMware不能直接使用这些普通的异步网络优化，而必须让
相关状态更新经过hypervisor的可记录路径：记录更新内容及其发生位置，然后在
backup VM replay时于相同位置应用同样的更新。这个问题和磁盘read completion
中通过DMA写入guest memory、随后投递interrupt的问题本质类似，都是外部设备
异步改变VM状态所带来的non-determinism。}

消除网络设备的异步更新，再加上section 2.2所述的发送包latency，%
给网络性能带来了一些挑战。在运行FT时，我们采用两种方法来提升VM网络性能。%
首先，我们实现了聚合优化以减少VM trap和interrupt。%
当VM以足够的比特率流式传输数据时，hypervisor可以对一组包只执行一次发%
送trap；在最佳情况下甚至可以零trap，因为它可以在接收新包的过程中发送这些%
包。同样，hypervisor也可以通过只为一组包投递一次interrupt，%
减少传入包触发给VM的interrupt数量。%

我们的第二个网络性能优化，是减少发送包的latency。如前所述，%
hypervisor必须delay所有发送包，直到它从backup VM收到相%
应log entry的acknowledgment。减少发送latency的关键，%
是减少向backup VM发送日志消息并获得acknowledgment所需的时间%
。我们在这方面的主要优化，是确保发送和接收log entry与%
acknowledgment都可以在没有任何thread context %
switch的情况下完成。VMware vSphere hypervisor允许向%
TCP栈注册函数，使其在收到TCP数据时从deferred-execution上下%
文(类似Linux中的tasklet)中被调用。这使我们能够在没有任何%
thread context switch的情况下，快速处理backup VM上收到的%
任何日志消息，以及primary VM上收到的任何acknowledgment。%
此外，当primary VM将一个待发送包入队时，我们会通过调度一个%
deferred-execution上下文来执行刷新，从而强制立即刷新相关联的输出%
log entry(如section 2.2所述)。
